\section{Executable Runtime Specification}
\label{app:runtime-specification}

This appendix specifies the control path used by the prototype. It is not a claim
that the following algorithms have been benchmarked on AppWorld or $\tau^2$-bench.

\paragraph{Main loop.} For task $i$, initialize a validated contract and durable
aggregate. While no budget is exhausted, build a bounded Context Pack; obtain a
model proposal; normalize each proposal into Action IR; validate schema, policy,
dependencies, and effect preflight; execute the allowed schedule; reduce tool
observations; verify state-changing effects; classify any failure; apply at most the
configured typed recovery; and assess progress and completion. Persist state and
trace records before another model call or offline evaluation. A rejected or
deferred call receives a structured result and is never silently discarded.

\paragraph{Contract Builder.} Given instruction $u$, active schemas $\Sigma$, and
policy $P$, produce candidate $Q$. Validate $Q$ against
\code{agentgate-core/src/agentgate\_core/schemas/v1/task-contract.schema.json};
canonicalize identifiers; reject dependency cycles and unsupported verifier
recipes; and place missing fields in \code{unresolved}. A deterministic normalizer
precedes at most one schema-constrained repair. The reducer records the accepted
contract version. The model cannot mark a goal completed in this procedure.

\paragraph{Completion Gate.} For every required subgoal, require
\code{completed\_verified}. For every required evidence and effect identifier,
require a fresh verified record. Reject on unresolved failure, pending confirmation
or approval, active evidence conflict, an \emph{active unresolved} unintended side
effect, or duplicate irreversible effect. An unexpected effect with status
\code{observed}, \code{rollback\_pending}, or \code{unresolved} blocks completion;
\code{rolled\_back} and \code{policy\_accepted} do not. Return every blocker with
its type, reference, and recommended
phase; accept only if the blocker set is empty.

\paragraph{Effect fingerprint.} For a state-changing Action IR $a$, canonicalize
the tool name, operation, resource type and identifier, and only the schema-marked
critical arguments. Serialize the canonical object with sorted keys and compute
\code{sha256}. The ledger uses this value as the idempotency key. A matching effect
that is verified yields \code{already\_applied}; an in-flight or unknown effect
yields \code{verify\_first}; a failed effect yields \code{recovery\_required}; and
a rolled-back effect may be reserved again. Effect reservation and aggregate
persistence use an optimistic compare-and-set revision. A competing reservation
that loses the idempotency or persistence race is not dispatched; it receives a
state-refresh/replan result. The implementation state machine is
\code{ABSENT $\rightarrow$ RESERVED $\rightarrow$ IN\_FLIGHT $\rightarrow$
APPLIED\_UNVERIFIED $\rightarrow$ VERIFIED}. A restored \code{IN\_FLIGHT} effect
becomes \code{UNKNOWN} and must follow \code{VERIFY\_FIRST}; \code{FAILED} follows
typed recovery; \code{VERIFIED $\rightarrow$ ROLLED\_BACK $\rightarrow$ RESERVED}
is the only replay path. This protocol is atomic at the persisted task-aggregate
revision, not a claim of an independently normalized SQL effect table.

\begin{center}
\small
\begin{tabularx}{\linewidth}{@{}p{0.23\linewidth}p{0.23\linewidth}YY@{}}
\toprule
Invariant & Enforcement location & Assumption & Local counterexample test \\
\midrule
Verified non-idempotent effect is not re-dispatched & Effect Ledger reservation and CAS persistence & every effectful tool is mapped to an effect & duplicate key; two SQLite store instances race \\
Completion has verified intended effects and no active unexpected effect & Completion Gate and verifier results & contract and verifier recipes are complete & missing effect; observed versus rolled-back/policy-accepted unexpected effect \\
Model cannot mutate authoritative state & typed State Reducer and store boundary & model process has no store capability & invalid direct transition and reducer transition tests \\
\bottomrule
\end{tabularx}
\captionof{table}{Runtime invariants, enforcement assumptions, and selected local
counterexample tests. This is executable test evidence, not a formal proof.}
\end{center}

\begin{center}
\small
\begin{tabularx}{0.94\linewidth}{@{}p{0.20\linewidth}YY@{}}
\toprule
Failure class & Typed recovery & Boundary \\
\midrule
Schema or parameter & deterministic normalization, then one constrained repair & no guessed fields \\
Entity & re-query and bind a unique identifier & ambiguous entities block \\
Transient tool & bounded retry after backoff & retry budget is finite \\
Environment mismatch & read back and reconcile expected effect & no blind write replay \\
Duplicate action & verify existing effect & non-idempotent writes do not replay \\
No progress or context loss & local replanning with state references & safe stop on repeated stagnation \\
Permission or policy & report blocker / request confirmation & runtime never bypasses policy \\
\bottomrule
\end{tabularx}
\captionof{table}{Failure-to-recovery mapping. Recovery success means the local
exit predicate is repaired; it does not imply evaluator-confirmed task success.}
\end{center}

\begin{center}
\small
\begin{tabular}{@{}lrp{0.53\linewidth}@{}}
\toprule
Context component & Budget priority & Must retain \\
\midrule
Task contract and hard constraints & fixed & goal IDs, prohibitions, confirmations \\
Current state and blockers & fixed & entity IDs, dependencies, unresolved failures \\
Tool schemas & bounded & active tools and required completion tool \\
Recent interactions & bounded & latest tool result and recovery context \\
Evidence and effects & bounded & verified effects, idempotency keys, freshness \\
Older history & references only & trace IDs and deterministic summaries \\
\bottomrule
\end{tabular}
\captionof{table}{Context Budget policy. The future experiment must freeze numeric
allocations in its run manifest; this paper does not claim a benchmark-optimized
allocation.}
\end{center}

\paragraph{Contract example.} The following abbreviated object illustrates the
versioned Task Contract schema. It exposes goals/subgoals, dependencies, hard and
policy constraints, allowed and forbidden effects, confirmation requirements, and
evidence requirements. Unresolved fields are deliberately represented in Task
State rather than the immutable contract, as specified in \cref{sec:contract-induction}.
\begin{verbatim}
{"task_id":"example-1", "version":1,
 "original_instruction":"Create and verify one quarterly record.",
 "normalized_goal":"Create one verified quarterly record.",
 "subgoals":[{"subgoal_id":"create", "description":"Create record",
 "dependency_ids":[], "completion_condition_ids":["record-verified"],
 "required":true}],
 "constraints":[{"constraint_id":"no-duplicates",
 "description":"Do not create a duplicate record", "kind":"hard"},
 {"constraint_id":"policy", "description":"Use catalog tools only",
 "kind":"policy"}],
 "completion_conditions":[{"condition_id":"record-verified",
 "description":"Readback verifies the record", "required_effect_ids":["effect-42"]}],
 "allowed_effects":["create_record"], "forbidden_effects":["delete_record"],
 "required_confirmations":[]}
\end{verbatim}
